% --- Archon physics formalization source begin ---
% archon:physics
% archon:covers PhyXMiniProblems/problem_phyx_mini_0192.lean
% archon:source-report reports/phyx_mini/problem_phyx_mini_0192.source.json

\chapter{Physics problem phyx\_mini\_0192}
\label{ch:PhyXMiniProblems_problem_phyx_mini_0192}

\paragraph{Problem source.}
\#\# Physical scenario

A directional loudspeaker directs a sound wave of wavelength \textbackslash{}lambda at a wall (figure).

\#\# Question

At what distances from the wall could you stand and hear no sound at all?

\#\# Answer choices

- A: A: d = \textbackslash{}frac\{\textbackslash{}lambda\}\{4\}

- B: B: d = \textbackslash{}frac\{\textbackslash{}lambda\}\{2\}

- C: C: d = \textbackslash{}frac\{\textbackslash{}lambda\}\{8\}

- D: D: d = \textbackslash{}frac\{\textbackslash{}lambda\}\{3\}

\#\# Figure caption (auxiliary; use the image as primary evidence)

The image depicts a sound wave interference diagram with the following components:

- A speaker is located on the left side, emitting sound waves.
- There is a tube or waveguide that gradually fades in color from dark to light blue, illustrating the wave propagation.
- On the right side, there is a wall reflecting the waves.
- Letters "N" and "A" are marked at different points along the wave path:
  - The pattern alternates between "N" (nodes) and "A" (antinodes).
- Three horizontal lines with measurements indicate distances:
  - \textbackslash{}( \textbackslash{}lambda/4 \textbackslash{})
  - \textbackslash{}( 3\textbackslash{}lambda/4 \textbackslash{})
  - \textbackslash{}( 5\textbackslash{}lambda/4 \textbackslash{})
- These measured lines show the distances between the nodes and antinodes relative to the wavelength \textbackslash{}( \textbackslash{}lambda \textbackslash{}).

\#\# Dataset metadata

Resonance and Harmonics; Physical Model Grounding Reasoning, Spatial Relation Reasoning

\paragraph{Recorded answer/context.}
A: A: d = \textbackslash{}frac\{\textbackslash{}lambda\}\{4\}

\paragraph{Figure/image path.}
/root/proposal\_for\_physic/science-mango/phyx\_mini\_run/phyx\_data/test\_image/192.png

\paragraph{Formalization target.}
create a compiling Lean file with sorry bodies at `PhyXMiniProblems/problem\_phyx\_mini\_0192.lean`. The Lean declarations must preserve the physical quantities, dimensions or dimensional roles, figure labels, governing-law hypotheses, and final relation expressed by this problem.
Use Mathlib/Physlib names found through LeanExplore where available. If a physics API is missing, introduce faithful local abstractions rather than scalar placeholder aliases.

\begin{theorem}[Physics formalization target]
\label{thm:physics:phyx_mini_0192:target}
\lean{PhyXMiniProblems.ProblemPhyXMini0192.problem_phyx_mini_0192}
\uses{def:physics:phyx-mini-0192:phyxminiproblems-problemphyxmini0192-lengthinmeters, def:physics:phyx-mini-0192:phyxminiproblems-problemphyxmini0192-loudspeakerwallsetup, def:physics:phyx-mini-0192:phyxminiproblems-problemphyxmini0192-isonwavepath, def:physics:phyx-mini-0192:phyxminiproblems-problemphyxmini0192-matchesproblemandfigure, def:physics:phyx-mini-0192:phyxminiproblems-problemphyxmini0192-hasphysicalacousticparameters, def:physics:phyx-mini-0192:phyxminiproblems-problemphyxmini0192-satisfiesrigidwallstandingwavepressurelaw, def:physics:phyx-mini-0192:phyxminiproblems-problemphyxmini0192-issilentstandinglocation, def:physics:phyx-mini-0192:phyxminiproblems-problemphyxmini0192-isphysicallycorrectchoice, def:physics:phyx-mini-0192:phyxminiproblems-problemphyxmini0192-recordeddatasetanswer}
The assigned autoformalize agent should translate this physics problem into Lean declarations in the covered file, with theorem and lemma proof bodies written as `by sorry`.
\end{theorem}
\begin{proof}
This is an autoformalization task, not a proof task. Produce statements that can later be proved by the physics prover without weakening the physical model.
\end{proof}

% --- Archon declaration topology begin ---
\section{Declaration topology}
\begin{definition}[Length Quantity]
  \label{def:physics:phyx-mini-0192:phyxminiproblems-problemphyxmini0192-lengthquantity}
  \lean{PhyXMiniProblems.ProblemPhyXMini0192.LengthQuantity}
  A nonnegative physical length, independent of the unit used to read it
\end{definition}
\begin{proof}
  This is the declared model definition of Length Quantity.
\end{proof}
\begin{definition}[Acoustic Pressure Amplitude]
  \label{def:physics:phyx-mini-0192:phyxminiproblems-problemphyxmini0192-acousticpressureamplitude}
  \lean{PhyXMiniProblems.ProblemPhyXMini0192.AcousticPressureAmplitude}
  A signed acoustic excess-pressure amplitude
\end{definition}
\begin{proof}
  This is the declared model definition of Acoustic Pressure Amplitude.
\end{proof}
\begin{definition}[length Readout]
  \label{def:physics:phyx-mini-0192:phyxminiproblems-problemphyxmini0192-lengthreadout}
  \lean{PhyXMiniProblems.ProblemPhyXMini0192.lengthReadout}
  \uses{def:physics:phyx-mini-0192:phyxminiproblems-problemphyxmini0192-lengthquantity}
  Read a physical length as a scalar in a selected length unit
\end{definition}
\begin{proof}
  This is the declared model definition of length Readout.
\end{proof}
\begin{definition}[length In Meters]
  \label{def:physics:phyx-mini-0192:phyxminiproblems-problemphyxmini0192-lengthinmeters}
  \lean{PhyXMiniProblems.ProblemPhyXMini0192.lengthInMeters}
  \uses{def:physics:phyx-mini-0192:phyxminiproblems-problemphyxmini0192-lengthquantity, def:physics:phyx-mini-0192:phyxminiproblems-problemphyxmini0192-lengthreadout}
  Meter readout of a physical distance or wavelength
\end{definition}
\begin{proof}
  This is the declared model definition of length In Meters.
\end{proof}
\begin{definition}[pressure Amplitude In Pascals]
  \label{def:physics:phyx-mini-0192:phyxminiproblems-problemphyxmini0192-pressureamplitudeinpascals}
  \lean{PhyXMiniProblems.ProblemPhyXMini0192.pressureAmplitudeInPascals}
  \uses{def:physics:phyx-mini-0192:phyxminiproblems-problemphyxmini0192-acousticpressureamplitude}
  Pascal readout of a signed acoustic excess-pressure amplitude
\end{definition}
\begin{proof}
  This is the declared model definition of pressure Amplitude In Pascals.
\end{proof}
\begin{definition}[Acoustic Source Kind]
  \label{def:physics:phyx-mini-0192:phyxminiproblems-problemphyxmini0192-acousticsourcekind}
  \lean{PhyXMiniProblems.ProblemPhyXMini0192.AcousticSourceKind}
  Kind of acoustic source on the left of the primary image
\end{definition}
\begin{proof}
  This is the declared model definition of Acoustic Source Kind.
\end{proof}
\begin{definition}[Acoustic Boundary Condition]
  \label{def:physics:phyx-mini-0192:phyxminiproblems-problemphyxmini0192-acousticboundarycondition}
  \lean{PhyXMiniProblems.ProblemPhyXMini0192.AcousticBoundaryCondition}
  Boundary behavior of the vertical wall on the right of the image
\end{definition}
\begin{proof}
  This is the declared model definition of Acoustic Boundary Condition.
\end{proof}
\begin{definition}[Acoustic Wave Kind]
  \label{def:physics:phyx-mini-0192:phyxminiproblems-problemphyxmini0192-acousticwavekind}
  \lean{PhyXMiniProblems.ProblemPhyXMini0192.AcousticWaveKind}
  Kind of wave established between the source and reflecting wall
\end{definition}
\begin{proof}
  This is the declared model definition of Acoustic Wave Kind.
\end{proof}
\begin{definition}[Displacement Role]
  \label{def:physics:phyx-mini-0192:phyxminiproblems-problemphyxmini0192-displacementrole}
  \lean{PhyXMiniProblems.ProblemPhyXMini0192.DisplacementRole}
  Displacement-amplitude labels denoted by N and A in the figure
\end{definition}
\begin{proof}
  This is the declared model definition of Displacement Role.
\end{proof}
\begin{definition}[Figure Point]
  \label{def:physics:phyx-mini-0192:phyxminiproblems-problemphyxmini0192-figurepoint}
  \lean{PhyXMiniProblems.ProblemPhyXMini0192.FigurePoint}
  Named axial locations in increasing distance from the wall. The last point is the additional N visible just to the right of the loudspeaker beam
\end{definition}
\begin{proof}
  This is the declared model definition of Figure Point.
\end{proof}
\begin{definition}[Standing Wave Figure]
  \label{def:physics:phyx-mini-0192:phyxminiproblems-problemphyxmini0192-standingwavefigure}
  \lean{PhyXMiniProblems.ProblemPhyXMini0192.StandingWaveFigure}
  \uses{def:physics:phyx-mini-0192:phyxminiproblems-problemphyxmini0192-lengthquantity, def:physics:phyx-mini-0192:phyxminiproblems-problemphyxmini0192-displacementrole, def:physics:phyx-mini-0192:phyxminiproblems-problemphyxmini0192-figurepoint}
  Positions and N/A annotations read from the supplied standing-wave image
\end{definition}
\begin{proof}
  This is the declared model definition of Standing Wave Figure.
\end{proof}
\begin{definition}[Loudspeaker Wall Setup]
  \label{def:physics:phyx-mini-0192:phyxminiproblems-problemphyxmini0192-loudspeakerwallsetup}
  \lean{PhyXMiniProblems.ProblemPhyXMini0192.LoudspeakerWallSetup}
  \uses{def:physics:phyx-mini-0192:phyxminiproblems-problemphyxmini0192-lengthquantity, def:physics:phyx-mini-0192:phyxminiproblems-problemphyxmini0192-acousticpressureamplitude, def:physics:phyx-mini-0192:phyxminiproblems-problemphyxmini0192-acousticsourcekind, def:physics:phyx-mini-0192:phyxminiproblems-problemphyxmini0192-acousticboundarycondition, def:physics:phyx-mini-0192:phyxminiproblems-problemphyxmini0192-acousticwavekind, def:physics:phyx-mini-0192:phyxminiproblems-problemphyxmini0192-standingwavefigure}
  The physical standing-wave experiment. pressureAmplitudeAt is the unknown spatial pressure-amplitude field of the resulting wave, and wallPressureAmplitude is its nonzero scale at the reflecting wall. Neither field is initialized with the requested quiet-location relation
\end{definition}
\begin{proof}
  This is the declared model definition of Loudspeaker Wall Setup.
\end{proof}
\begin{definition}[Is On Wave Path]
  \label{def:physics:phyx-mini-0192:phyxminiproblems-problemphyxmini0192-isonwavepath}
  \lean{PhyXMiniProblems.ProblemPhyXMini0192.IsOnWavePath}
  \uses{def:physics:phyx-mini-0192:phyxminiproblems-problemphyxmini0192-lengthquantity, def:physics:phyx-mini-0192:phyxminiproblems-problemphyxmini0192-lengthinmeters, def:physics:phyx-mini-0192:phyxminiproblems-problemphyxmini0192-loudspeakerwallsetup}
  A distance lies on the modeled waveguide segment between wall and speaker
\end{definition}
\begin{proof}
  This is the declared model definition of Is On Wave Path.
\end{proof}
\begin{definition}[Matches Problem And Figure]
  \label{def:physics:phyx-mini-0192:phyxminiproblems-problemphyxmini0192-matchesproblemandfigure}
  \lean{PhyXMiniProblems.ProblemPhyXMini0192.MatchesProblemAndFigure}
  \uses{def:physics:phyx-mini-0192:phyxminiproblems-problemphyxmini0192-lengthinmeters, def:physics:phyx-mini-0192:phyxminiproblems-problemphyxmini0192-loudspeakerwallsetup, def:physics:phyx-mini-0192:phyxminiproblems-problemphyxmini0192-isonwavepath}
  The prose setup and all discrete or measured information taken from the primary image. The figure displays a displacement node at the wall and then alternating antinodes and nodes at quarter-wavelength intervals. Its three dimension arrows terminate at the antinodes lambda/4, 3 lambda/4, and 5 lambda/4 from the wall. These are geometric and displacement-role readouts only. No field asserts zero pressure, silence, or correctness of an answer choice
\end{definition}
\begin{proof}
  This is the declared model definition of Matches Problem And Figure.
\end{proof}
\begin{definition}[Has Physical Acoustic Parameters]
  \label{def:physics:phyx-mini-0192:phyxminiproblems-problemphyxmini0192-hasphysicalacousticparameters}
  \lean{PhyXMiniProblems.ProblemPhyXMini0192.HasPhysicalAcousticParameters}
  \uses{def:physics:phyx-mini-0192:phyxminiproblems-problemphyxmini0192-lengthinmeters, def:physics:phyx-mini-0192:phyxminiproblems-problemphyxmini0192-pressureamplitudeinpascals, def:physics:phyx-mini-0192:phyxminiproblems-problemphyxmini0192-loudspeakerwallsetup}
  Positivity and nondegeneracy conditions for an audible incident wave
\end{definition}
\begin{proof}
  This is the declared model definition of Has Physical Acoustic Parameters.
\end{proof}
\begin{definition}[Satisfies Rigid Wall Standing Wave Pressure Law]
  \label{def:physics:phyx-mini-0192:phyxminiproblems-problemphyxmini0192-satisfiesrigidwallstandingwavepressurelaw}
  \lean{PhyXMiniProblems.ProblemPhyXMini0192.SatisfiesRigidWallStandingWavePressureLaw}
  \uses{def:physics:phyx-mini-0192:phyxminiproblems-problemphyxmini0192-lengthinmeters, def:physics:phyx-mini-0192:phyxminiproblems-problemphyxmini0192-pressureamplitudeinpascals, def:physics:phyx-mini-0192:phyxminiproblems-problemphyxmini0192-loudspeakerwallsetup, def:physics:phyx-mini-0192:phyxminiproblems-problemphyxmini0192-isonwavepath}
  The governing pressure profile for a one-dimensional wave reflected from a rigid wall. Taking distance x outward from the wall, the excess-pressure amplitude is P\_wall cos(2 pi x / lambda). Thus the wall itself is a pressure antinode even though it is a displacement node. This is a general standing-wave law. It does not list the zeros of the cosine and does not state any requested distance or answer choice
\end{definition}
\begin{proof}
  This is the declared model definition of Satisfies Rigid Wall Standing Wave Pressure Law.
\end{proof}
\begin{definition}[Is Silent Standing Location]
  \label{def:physics:phyx-mini-0192:phyxminiproblems-problemphyxmini0192-issilentstandinglocation}
  \lean{PhyXMiniProblems.ProblemPhyXMini0192.IsSilentStandingLocation}
  \uses{def:physics:phyx-mini-0192:phyxminiproblems-problemphyxmini0192-lengthquantity, def:physics:phyx-mini-0192:phyxminiproblems-problemphyxmini0192-pressureamplitudeinpascals, def:physics:phyx-mini-0192:phyxminiproblems-problemphyxmini0192-loudspeakerwallsetup, def:physics:phyx-mini-0192:phyxminiproblems-problemphyxmini0192-isonwavepath}
  A listener hears no sound from the modeled mode at a location on the wave path exactly when the acoustic pressure amplitude there vanishes
\end{definition}
\begin{proof}
  This is the declared model definition of Is Silent Standing Location.
\end{proof}
\begin{definition}[Answer Choice]
  \label{def:physics:phyx-mini-0192:phyxminiproblems-problemphyxmini0192-answerchoice}
  \lean{PhyXMiniProblems.ProblemPhyXMini0192.AnswerChoice}
  Labels of the four distance choices printed with the problem
\end{definition}
\begin{proof}
  This is the declared model definition of Answer Choice.
\end{proof}
\begin{definition}[displayed Wavelength Ratio]
  \label{def:physics:phyx-mini-0192:phyxminiproblems-problemphyxmini0192-displayedwavelengthratio}
  \lean{PhyXMiniProblems.ProblemPhyXMini0192.displayedWavelengthRatio}
  \uses{def:physics:phyx-mini-0192:phyxminiproblems-problemphyxmini0192-answerchoice}
  Dimensionless coefficient of lambda displayed beside an answer label
\end{definition}
\begin{proof}
  This is the declared model definition of displayed Wavelength Ratio.
\end{proof}
\begin{definition}[Matches Displayed Distance]
  \label{def:physics:phyx-mini-0192:phyxminiproblems-problemphyxmini0192-matchesdisplayeddistance}
  \lean{PhyXMiniProblems.ProblemPhyXMini0192.MatchesDisplayedDistance}
  \uses{def:physics:phyx-mini-0192:phyxminiproblems-problemphyxmini0192-lengthquantity, def:physics:phyx-mini-0192:phyxminiproblems-problemphyxmini0192-lengthinmeters, def:physics:phyx-mini-0192:phyxminiproblems-problemphyxmini0192-loudspeakerwallsetup, def:physics:phyx-mini-0192:phyxminiproblems-problemphyxmini0192-answerchoice, def:physics:phyx-mini-0192:phyxminiproblems-problemphyxmini0192-displayedwavelengthratio}
  A physical distance has the wavelength ratio printed beside a choice
\end{definition}
\begin{proof}
  This is the declared model definition of Matches Displayed Distance.
\end{proof}
\begin{definition}[Is Physically Correct Choice]
  \label{def:physics:phyx-mini-0192:phyxminiproblems-problemphyxmini0192-isphysicallycorrectchoice}
  \lean{PhyXMiniProblems.ProblemPhyXMini0192.IsPhysicallyCorrectChoice}
  \uses{def:physics:phyx-mini-0192:phyxminiproblems-problemphyxmini0192-loudspeakerwallsetup, def:physics:phyx-mini-0192:phyxminiproblems-problemphyxmini0192-issilentstandinglocation, def:physics:phyx-mini-0192:phyxminiproblems-problemphyxmini0192-answerchoice, def:physics:phyx-mini-0192:phyxminiproblems-problemphyxmini0192-matchesdisplayeddistance}
  A choice is physically possible when a silent point has its displayed ratio
\end{definition}
\begin{proof}
  This is the declared model definition of Is Physically Correct Choice.
\end{proof}
\begin{definition}[recorded Dataset Answer]
  \label{def:physics:phyx-mini-0192:phyxminiproblems-problemphyxmini0192-recordeddatasetanswer}
  \lean{PhyXMiniProblems.ProblemPhyXMini0192.recordedDatasetAnswer}
  \uses{def:physics:phyx-mini-0192:phyxminiproblems-problemphyxmini0192-answerchoice}
  Answer label recorded in the supplied dataset metadata
\end{definition}
\begin{proof}
  This is the declared model definition of recorded Dataset Answer.
\end{proof}
\begin{lemma}[silent iff cosine phase zero]
  \label{lem:physics:phyx-mini-0192:phyxminiproblems-problemphyxmini0192-silent-iff-cosine-phase-zero}
  \lean{PhyXMiniProblems.ProblemPhyXMini0192.silent_iff_cosine_phase_zero}
  \uses{def:physics:phyx-mini-0192:phyxminiproblems-problemphyxmini0192-lengthquantity, def:physics:phyx-mini-0192:phyxminiproblems-problemphyxmini0192-lengthinmeters, def:physics:phyx-mini-0192:phyxminiproblems-problemphyxmini0192-loudspeakerwallsetup, def:physics:phyx-mini-0192:phyxminiproblems-problemphyxmini0192-isonwavepath, def:physics:phyx-mini-0192:phyxminiproblems-problemphyxmini0192-hasphysicalacousticparameters, def:physics:phyx-mini-0192:phyxminiproblems-problemphyxmini0192-satisfiesrigidwallstandingwavepressurelaw, def:physics:phyx-mini-0192:phyxminiproblems-problemphyxmini0192-issilentstandinglocation}
  Under the rigid-wall pressure law, silence is equivalent to the vanishing of the cosine phase factor. This intermediate result uses nonzero incident amplitude but does not solve the cosine-zero equation
\end{lemma}
\begin{proof}
  The conclusion follows from the governing relations and figure readouts named in the statement of silent iff cosine phase zero.
\end{proof}
\begin{lemma}[silent locations are odd quarter wavelengths]
  \label{lem:physics:phyx-mini-0192:phyxminiproblems-problemphyxmini0192-silent-locations-are-odd-quarter-wavelengths}
  \lean{PhyXMiniProblems.ProblemPhyXMini0192.silent_locations_are_odd_quarter_wavelengths}
  \uses{def:physics:phyx-mini-0192:phyxminiproblems-problemphyxmini0192-lengthinmeters, def:physics:phyx-mini-0192:phyxminiproblems-problemphyxmini0192-loudspeakerwallsetup, def:physics:phyx-mini-0192:phyxminiproblems-problemphyxmini0192-isonwavepath, def:physics:phyx-mini-0192:phyxminiproblems-problemphyxmini0192-hasphysicalacousticparameters, def:physics:phyx-mini-0192:phyxminiproblems-problemphyxmini0192-satisfiesrigidwallstandingwavepressurelaw, def:physics:phyx-mini-0192:phyxminiproblems-problemphyxmini0192-issilentstandinglocation}
  The complete family of quiet listening locations is the set of nonnegative odd quarter-wavelengths which still lie between the wall and loudspeaker. This conclusion is derived from the pressure law; it is not a premise
\end{lemma}
\begin{proof}
  The conclusion follows from the governing relations and figure readouts named in the statement of silent locations are odd quarter wavelengths.
\end{proof}
\begin{lemma}[first pictured antinode is silent]
  \label{lem:physics:phyx-mini-0192:phyxminiproblems-problemphyxmini0192-first-pictured-antinode-is-silent}
  \lean{PhyXMiniProblems.ProblemPhyXMini0192.first_pictured_antinode_is_silent}
  \uses{def:physics:phyx-mini-0192:phyxminiproblems-problemphyxmini0192-loudspeakerwallsetup, def:physics:phyx-mini-0192:phyxminiproblems-problemphyxmini0192-matchesproblemandfigure, def:physics:phyx-mini-0192:phyxminiproblems-problemphyxmini0192-hasphysicalacousticparameters, def:physics:phyx-mini-0192:phyxminiproblems-problemphyxmini0192-satisfiesrigidwallstandingwavepressurelaw, def:physics:phyx-mini-0192:phyxminiproblems-problemphyxmini0192-issilentstandinglocation}
  The first A in the image, marked at lambda/4 from the wall, is a quiet pressure node. This connects the general pressure law to the particular figure readout without treating the displayed A displacement label itself as an assumption of silence
\end{lemma}
\begin{proof}
  The conclusion follows from the governing relations and figure readouts named in the statement of first pictured antinode is silent.
\end{proof}
% --- Archon declaration topology end ---
% --- Archon physics formalization source end ---
